\section{Wprowadzenie do analizy}
\lipsum[1] % Zwykły tekst akapitu.

% Użycie środowiska 'verbetering'
\begin{remark}
To jest ważna uwaga dotycząca zbieżności szeregów.
Pamiętaj, aby zawsze sprawdzać warunek konieczny. Ten boks ma stały tytuł "Opmerking" i ciemnozieloną ramkę.
\end{remark}

\[
	\mfaktor{A}{B}
.\] 

\lipsum[2]
A teraz przejdziemy do definicji. 
% Użycie środowiska 'noot' z własnym tytułem
\begin{note}{Definicja: Ciągłość funkcji}
Mówimy, że funkcja $f(x)$ jest ciągła w punkcie $x_0$, jeśli spełnione są trzy warunki:
\begin{enumerate}
    \item Istnieje wartość funkcji $f(x_0)$.
    \item Istnieje granica $\lim_{x \to x_0} f(x)$.
    \item Granica ta jest równa wartości funkcji: $\lim_{x \to x_0} f(x) = f(x_0)$.
\end{enumerate}
Ten boks ma szarą ramkę, a jego tytuł został podany jako argument.
\end{note}

\lipsum[3]
\begin{definition}[Podgrupa grupy]	

	Podgrupą grupy G nazywamy taki podzbiór $H\subseteq G$, który spełnia następujące warunki:
\begin{itemize}
	\item Zamkniętość na mnożenie-  $\left( \forall x,y \in H \right) x\cdot y \in H$
	\item Zamkniętość na branie elementu odwrotnego- $\left( \forall x \in H \right)x^{-1} \in H$
	\item Element neutralny jest zawarty- $\mathbf{1}\in H$
\end{itemize}
\end{definition}

Zapis  $H\leq G$ będzie oznaczać, że H jest podgrupą grupy G.

Zauważmy, że $\mathbf{1} = \{1\} \le G$. Taką podgrupę będziemy nazywać podgrupą trywialną. Cała grupa G także jest swoją podgrupą, tzn. $G\le G$.

Przykłady:
\begin{enumerate}
	\item Tutaj przykłady
	\item Dalsze przykłady
\end{enumerate}

\begin{definition}{Homomorfizm grup.}

	Niech G, H- grupy.
	Przekształcenie $\varphi: G \rightarrow H$ nazywamy homomorfizmem grup wtedy i tylko wtedy gdy $\left( \forall g_1, g_2 \in G \right) \varphi(g_1\cdot g_2) = \varphi(g_1)\cdot \varphi(g_2)$
\end{definition}

(Cwiczenie) Homomorfizm $\varphi$ przeprowadza element neutralny na element neutralny, a element odwrotny do g na element odwrotny $\varphi(g)$.
Istnieją rożne szczególne rodzaje homomorfizmów:

\begin{itemize}
	\item Izomorfizm. Jest to taki homomorfizm $\varphi: G \rightarrow H$, dla którego istnieje homomorfizm $\psi : H \rightarrow G$, taki że $\varphi\psi = \text{id}_{H}$ oraz $\psi\varphi = \text{id}_{G}$.
		\begin{remark}

			Homomorfizm grup jest izomorfizmem wtedy i tylko wtedy kiedy jest homomorfizmem i bijekcją zbiorów.
		\end{remark}
		Grupy izomorficzne będziemy uważać za takie same.
	\item coasa
\end{itemize}
\begin{proof}

	DOWo
	\lipsum[1]
\end{proof}


\lipsum[2]


\begin{theorem}[Lagrange'a]
    Jeśli G jest grupą skończoną, a H jej podgrupą, to $|G| = |H| \cdot [G:H]$.
\end{theorem}
\begin{proof}
	Z lematu o warstwach wiemy, że lewe warstwy H tworzą partycję\marginnote{Czyli jest relacją równoważności} zbioru G. Oznacza to, że G jest rozłączną sumą swoich warstw. Wszystkich warstw jest $[G:H]$. Z tego samego lematu wiemy, że każda warstwa ma moc $|H|$. Zatem sumując moce wszystkich warstw, otrzymujemy $|G| = \underbrace{|H| + |H| + \dots + |H|}_{[G:H] \text{ razy}} = [G:H] \cdot |H|$.
\end{proof}

\lipsum[2]

\begin{eg}
    Jeśli G jest grupą skończoną, a H jej podgrupą, to $|G| = |H| \cdot [G:H]$.
\end{eg}

\lipsum[1]

\begin{note}[Notatka]

	\lipsum[1]
\end{note}
\lipsum[2]
\begin{remark}[Inna notatka]
	\lipsum[1]	
\end{remark}
\lipsum[3]
\begin{corollary}[Wazny wniosek]
	\lipsum[1]	
\end{corollary}
\begin{ex}

	Udowodnić powyższe stwierdzenia. Następnie zrobić notatkę.
\end{ex}

